% !TeX root = ../drafts/01-introduction.tex
\section{Introduction}\label{sec:intro}

Ethereum's consensus layer currently uses BLS signatures, its execution layer ECDSA, and its data availability layer KZG. All three are based on elliptic curves, which do not survive a quantum computer~\cite{shor}. Hash-based cryptography, both signatures and snarks (the hash-based ones are sometimes called STARKs~\cite{stark}), is a promising candidate for Ethereum's post-quantum transition:
\begin{itemize}[itemsep=1pt,topsep=3pt]
\item BLS $\to$ XMSS~\cite{rfc8391} (using each Ethereum slot as the nonce, statefulness is no longer an issue when attesting: double signing is already a slashable offense). But XMSS has no native aggregation, unlike BLS, so we need a snark to handle it.
\item ECDSA $\to$ SPHINCS+~\cite{sphincs}. Drawback: dozens of times bigger. Again, snark-based aggregation is a natural solution to mitigate signature size increase.
\item KZG $\to$ Reed--Solomon encode each blob, and provide the Merkle root of the codeword. A snark guarantees that the root commits to a correct encoding.
\end{itemize}

\vspace{4mm}

leanVM is a candidate for the snark machinery required above. It proves execution of native RV64IM programs compiled from Rust, using a bounded byte-addressed machine and proof-checked cryptographic ECALLs. The machine uses ordinary $64$-bit integer registers and mutable byte RAM; binary fields belong to the proof system, not to its instruction-word semantics. The proof system retains WHIR, the binary-field arithmetization inspired by \cite{DP24}, and Flock \cite{Flock26} for BLAKE2s \cite{blake2}. It is commonly called a zkVM, although zero-knowledge is not currently provided.
